%  LaTeX support: latex@mdpi.com 
%  For support, please attach all files needed for compiling as well as the log file, and specify your operating system, LaTeX version, and LaTeX editor.

%=================================================================
\documentclass[materials,article,submit,pdftex,moreauthors]{Definitions/mdpi} 
% For posting an early version of this manuscript as a preprint, you may use "preprints" as the journal and change "submit" to "accept". The document class line would be, e.g., \documentclass[preprints,article,accept,moreauthors,pdftex]{mdpi}. This is especially recommended for submission to arXiv, where line numbers should be removed before posting. For preprints.org, the editorial staff will make this change immediately prior to posting.

%=================================================================
% MDPI internal commands
\firstpage{1} 
\makeatletter 
\setcounter{page}{\@firstpage} 
\makeatother
\pubvolume{1}
\issuenum{1}
\articlenumber{0}
\pubyear{2022}
\copyrightyear{2022}
%\externaleditor{Academic Editor: Firstname Lastname}
\datereceived{} 
%\daterevised{} % Only for the journal Acoustics
\dateaccepted{} 
\datepublished{} 
%\datecorrected{} % Corrected papers include a "Corrected: XXX" date in the original paper.
%\dateretracted{} % Corrected papers include a "Retracted: XXX" date in the original paper.
\hreflink{https://doi.org/} % If needed use \linebreak
%\doinum{}

%\bibstyle{numeric}
\usepackage[super]{nth}

%------------------------------------------------------------------
% The following line should be uncommented if the LaTeX file is uploaded to arXiv.org
%\pdfoutput=1

%=================================================================
% Add packages and commands here. The following packages are loaded in our class file: fontenc, inputenc, calc, indentfirst, fancyhdr, graphicx, epstopdf, lastpage, ifthen, lineno, float, amsmath, setspace, enumitem, mathpazo, booktabs, titlesec, etoolbox, tabto, xcolor, soul, multirow, microtype, tikz, totcount, changepage, attrib, upgreek, cleveref, amsthm, hyphenat, natbib, hyperref, footmisc, url, geometry, newfloat, caption

%=================================================================
%% Please use the following mathematics environments: Theorem, Lemma, Corollary, Proposition, Characterization, Property, Problem, Example, ExamplesandDefinitions, Hypothesis, Remark, Definition, Notation, Assumption
%% For proofs, please use the proof environment (the amsthm package is loaded by the MDPI class).

%=================================================================
% Full title of the paper (Capitalized)
\Title{Dynamics of strength gain in sandy soil amended with stabilized materials evaluated by elastic P-waves}

% MDPI internal command: Title for citation in the left column
\TitleCitation{Dynamics of strength gain in sandy soil amended with stabilized materials evaluated by elastic P-waves}

% Author Orchid ID: enter ID or remove command
\newcommand{\orcidauthorA}{0000-0002-0577-9936} % Add \orcidA{} behind the author's name
\newcommand{\orcidauthorB}{0000-0002-5759-1089} % Add \orcidA{} behind the author's name

% Authors, for the paper (add full first names)
\Author{Per Lindh $^{1,2}$\orcidA{} and Polina Lemenkova $^{3}\orcidB{}$*}

%\longauthorlist{yes}

% MDPI internal command: Authors, for metadata in PDF
\AuthorNames{Per Lindh, Polina Lemenkova}

% MDPI internal command: Authors, for citation in the left column
\AuthorCitation{Lindh, P.; Lemenkova, P.}
% If this is a Chicago style journal: Lastname, Firstname, Firstname Lastname, and Firstname Lastname.

% Affiliations / Addresses (Add [1] after \address if there is only one affiliation.)
\address{%
$^{1}$ \quad Swedish Transport Administration, Malm\"{o}, Sweden; per.a.lindh@gmail.com \\
$^{2}$ \quad Lund University, Lunds Tekniska H\"{o}gskola LTH (Faculty of Engineering), Department of Building and Environmental Technology, Division of Building Materials, Box 118, SE-221-00, Lund, Sweden.\\
$^{3}$ \quad Universit\'{e} Libre de Bruxelles, \'{E}cole polytechnique de Bruxelles (Brussels Faculty of Engineering), Laboratory of Image Synthesis and Analysis, Building L, Campus de Solbosch, ULB -- LISA CP165/57, Avenue Franklin D. Roosevelt 50, B-1050, Brussels, Belgium}

% Contact information of the corresponding author
\corres{Correspondence: polina.lemenkova@ulb.be; Tel.: +32471860459 (P. Lemenkova)}

% Current address and/or shared authorship
\firstnote{Current address: Affiliation 3.} 
\secondnote{These authors contributed equally to this work.}
% The commands \thirdnote{} till \eighthnote{} are available for further notes

%\simplesumm{} % Simple summary

%\conference{} % An extended version of a conference paper

% Abstract (Do not insert blank lines, i.e. \\) 
\abstract{Stabilization of poor subgrade soil intended for road construction is aimed at improving its engineering properties, commonly estimated as strength and stiffness. The improvement of these parameters is usually performed with binders, such as lime, slag, cement. The growth of strength and stiffness is strongly dependent on both time of curing and the effects from the stabilizing agents which can be used both as single binders and as blended mixtures. Using mixed binders as stabilising agents may have diverse effects on the properties of soil. Therefore, the experimental testing of the effects of blenders is required. This paper deals with experimental trials of five different binders for stabilization of clayey soil: cement, lime, slag (GGBFS) and two types of fly ash (energy and bio fly ash). The aim is to evaluate the effects from single and mixed binders on the gain of strength over time of curing. The optimal selection of binder combinations and binder amounts was tested using experimental planning with the aim at reducing the financial risks of the project, because large amounts of soil had to be stabilized in real project (several hundreds of tons of soil). Therefore, the optimization of binder blending was performed in a laboratory of SGI using modelling by mixture simplex lattice design with three binders in each case as independent variables. The performed methodology include soil treatment by 5 types of binders during 90 days, and geotechnical tests according to the SGI standards: strength test for the Unconfined Compressive Strength (UCS) by the SIS standard SS-EN 13286-41 and seismic tests on the stabilized samples using ICP Accelerometer for measuring P-wave velocity. The dynamics of soil behaviour stabilized by different binders for days 7, 14, 28 and 90 was statistically analysed and compared. The curves of strength gain in soil specimens have been compared using various ratios of binders (double and triple combinations of stabilizing agents). The comparative analysis was performed to assess the functional behavior of the P-wave velocity controlled by effects from individual and blended (binary and triple) binders and curing time of 90 days. Besides, variation in water ratio content in specimens before and upon the curing period of 90 days were statistically assessed. The development in P-wave velocity was compared both for individual and blended binders over the complete period of curing. Results indicate that strength growth in stabilised soil samples is nonlinear in both time and content of binders. The analysis of the effects from various agents demonstrated that slag contributes the most to the compressive strength development, followed by cement. Besides, we noted that blended binders perform better than single binders. Thus, the combination of lime and slag seems to have a notable effect on hardening process of soil, although effect on the final strength development. The combination of slag and energy ash contributes to the gain of strength, but has the least impact on compared to other mixtures.}

% Keywords
\keyword{civil engineering; stabilization; compressive strength; soil; cement; slag; fly ash; lime; seismic waves} 

% The fields PACS, MSC, and JEL may be left empty or commented out if not applicable
%\PACS{J0101}
%\MSC{}
%\JEL{}

%%%%%%%%%%%%%%%%%%%%%%%%%%%%%%%%%%%%%%%%%%
\begin{document}

\section{Introduction}

The goal of laboratory experiments on soil stabilization, given a context of road and highway construction in civil engineering, is to determine the compaction characteristics and the strength of the stabilised material. Stabilization of soil with traditional binders is used in a wide variety of engineering applications with the major aim at improving soil properties prior to construction works in road engineering (Källén et al. 2014, 2016; Abdullah and Al-Abdul Wahhab, 2015; Glossop and Wilson, 1953; Park et al. 2020; Lemenkov and Lemenkova, 2021a; Sudhakar et al. 2021; Nordmark et al. 2022). Cement and lime are the oldest, fundamental and most well-known traditional binders, continuously used in civil engineering works (Hilt and Davidson, 1960; Mitchell and El Jack, 1966; Kaniraj and Havanagi, 1999; Chen and Wang, 2006). Therefore, their applications for soil stabilization has been researched and reported for decades as a straightforward, versatile and traditional method of soil improvement (Markwick \& Keep, 1942; Bell, 1989; Dahlin et al. 1999; Aldaood et al. 2021; Abbey et al. 2021). 

Although traditional binders, such as cement or lime, remain the most commonly used additives in soil stabilization, the selection of binders because more complicated owing to recent challenges. These include negative environmental impacts from cement, the increased amount of the industrial by-products which should be utilized and the fabrication of novel stabilizing agents which have improved geotechnical properties in terms of engineering performance and better suited to the environmental challenges. For instance, such specific cases arise when geotechnical work are planned on the expansive or clayey soil that have specific properties (Nelson and Miller, 1992).  Blended binders, combined from several materials of various additives can response to such needs and requirements, since they are better suited to the different soil types, environmental conditions, thus considering specific setting of a real project.

In this regard, blended binders can be considered as an advanced method of soil stabilisation (Jafer, 2017; Jegandan et al. 2010). The experimental use of the alternative and blended binders allows for application of more effective approaches, both in economical and engineering aspects (Mavroulidou et al. 2021), which is valuable, due to the high cost of geotechnical works on road construction. Various attempts have been made to understand how the addition of of blended binders can enhance soil properties and increase its strength, and in which proportions binders should be taken (Verhasselt, 1990; Al-Swaidani et al. 2016; Sherwood, 1993; Fasihnikoutalab et al. 2020; Chaddock, 1996; Kogbara et al. 2011). Indeed, mixed binders have the potential to improve soil stabilization through more appropriate workflow of the experiments. For instance, the use of blended binders supports the effectiveness of clayey soil hardening over time of curing (Wang et al. 2015; Lindh \& Lemenkova, 2021a; Gandhi and Shukla, 2021). 

Specifically for Sweden, the performance of single and mixed binders was investigated using deep mixed method (Holm and Åhnberg, 1987; Åhnberg et al., 1995). The most well-known application of blended binders in Sweden is the depth stabilization technique, which follows the long-term tradition of soil treatment in the country (Bjerrum \& Flodin, 1960). Others include using blended binders for shallow soil stabilisation. Stabilization has also gained the increased use in infrastructure projects for treatment of the contaminated soil (Sörengård et al. 2021; Lindh \& Lemenkova, 2022b). Further, blended binders perform well compared to single binders, because such mixtures are better suitable on soils which with varied grain size (el Mahdi Safhi et al. 2022; Syed et al. 2022). Since the response of such soils varies to the different types of binders, blended mixes can be applied for testing soils with diverse grain content, ranging from a fine-grained clay to a coarse gravel. Moreover, the use of mixed binders can help make cemented soil structure estimates consistent with the in-situ conditions of real projects, such as road engineering. 

However, as the areas of applications increases along with the use of complex mixtures, so do the requirements for the optimization of soil. In the post-treatment of the contaminated soil, it is common for chemically active additives to be used to increase the effect of stabilization with regard to the reduction of the environmental risk. It is also common for cement or slag to be replaced by other binders, such as fly ash or other residual products, with aim to reduce costs of engineering works (Jafer et al. 2018; Lindh and Lemenkova, 2022a; Arulrajah et al. 2018). Besides, in response to the increased societal demands for environmentally sustainably solutions in engineering works, e.g., reduced CO2 emissions, the demands on the environmental impact from binders during the process of soil stabilization are also increasing. Overall, these requirements and new challenges in road construction industry lead to the need to gain access to methods that enable optimization of soil treatment that takes into account both technical performance, economics and environmental aspects to be considered in final project.

The UCS test is a straightforward and relatively easy method to evaluate the enhancement of strength and stiffness in stabilized soil over time. Because it well reflects both compaction properties and gained strength, it is one of the most frequently used and commonly acceptable methods to evaluate the engineering properties of the stabilized soil, as continuously reported in earlier papers (Christensen, 1969; Rodriguez et al. 1988; Eskisar, 2015; Lemenkov and Lemenkova, 2021b; Feng et al. 2021). The examples of using UCS tests are diverse and include, for instance, improved methods testing bearing capacity and settlement of cohesive soil in the foundation of building (Alshkane et al. 2020). Some more advanced methods include predictive modelling by programming techniques to estimate the UCS (Onyelowe et al. 2022), using artificial intelligence methods (Sihag et al. 2021; Liu et al. 2015; Tabarsa et al. 2021) and machine learning applications (Suthar, 2020; Furtney et al. 2022) to UCS estimation.

The alternative methods of evaluation of stiffness and strength parameters of soil stabilized by various percentages of binders can be performed using seismic tests measured velocities of elastic pressure waves (P-waves). The core principle of the evaluation of velocities of pressure wave consists in the physical theory of waves and elastic properties of soils as a porous media. Thus, measuring P-wave speed enables to estimate the level of stiffness and strength in a soil sample, as used in various geotechnical application (Gheibi and Hedayat, 2018; Jiang et al. 2017, 2021; Lindh and Lemenkova, 2021b; Lacidogna et al. 2019). The technical approach of this method consists in the propagation of elastic waves that penetrate the specimens of soil materials provided at different curing periods. A quantitative assessment of the P-wave velocity gives the value of soil strength and stiffness due to the correlation between these parameters, which is reported in relevant case studies (Ryden et al. 2006; Kurtulus et al. 2010; Lindh and Lemenkova, 2022b; Lin and Wang, 2020). Seismic methods applied to measuring properties of soil demonstrated superiority over common methods due to its non-destructive nature, which enables to measure samples many times, but requires complex tools and advanced methods of data processing (Shi et al. 2017; Sarro et al. 2021). 

This paper deals with experimental trials of five different binders for stabilization of clayey soil: cement, lime, GGBFS and two types of fly ash (energy and bio fly ash). The combinations of blended binders were determined using special techniques of simplex designs, aimed at optimising the proportions of the stabilizing agents. The methodology include soil treatment by blended binders during 90 days, and geotechnical tests according to the SGI standards: strength test for the UCS and seismic measurements (P-wave speed) of the stabilized specimens. The dynamics of soil behaviour stabilized by different binders was statistically analysed and compared on days 7, 14, 28 and 90th. The development in P-wave velocity was compared both for single and blended binders over the complete period of curing (90 days). The results are summarised in tables and visualized on plots in relevant chapters below. 

\section{Experimental workflow}

Workflow optimization using experimental design is an important part of geotechnical works, as it enables to minimize the number of tests. In view of the large amounts of material that should be tested (hundreds of tons), the optimization of work is necessary for economically more effective planning of the workload. To assess the impacts from individual and mixed binders, a simplex mixture design was chosen following the existing methodology (Myers and Montgomery, 1995). The role of trial planning consists of a systematic methodology aimed at finding the optimal and minimal necessary content of the tested binder products which should be added into the soil mixture. The traditional method to optimizing binder mixtures uses the approach which consists in one factor tested at a time. This strategy enables only one parameter to be changed, while keeping all the other parameters as constants. However, by this approach there is a risk of missing both positive and negative interactions between the constituent components of soil-binder mixture. Neither it allows a simultaneous optimization which could be based on the analysis of both technical and environmental requirements of stabilized soil.

The decision regarding which binder mixture is optimal is based on a series of different random trials where various factors can be varied in an unsystematic way. The mentioned above problems can be eliminated if the statistical trial planning is used to support the optimization of binder mixtures. Thus, the trial design ensures that there is a systematic evaluation of binder mixtures, which is achieved through simultaneous taking into account the influence of various factors on the final results (Montgomery, 1996). Specifically for soil testing, in this case we can measure the effects of various combination of binders on the evolution of soil strength over time of curing. By working systematically, the possibility of making a good optimization of blended mixtures of binders is increased with regard to both technical performance, economic gain and environmental impact.  

Laboratory tests on binder proportions and combinations often entails higher costs, which is caused due to many trials in the initial stage. Using this method, high total cost of geotechnical works on soil stabilization intended for road construction can be significantly reduced. Since both the choice of binder components and binder quantities largely affect project economy, the price of works can be reduced in relation to the needs of real project. Therefore, to optimize binder mixtures we used modelling trial design techniques. The approach can be divided into a mixing part and a process control part. Mixture optimization included the mutual proportions of different binders in the final mixture, that is, the ratios of blended binders (Lindh, 2001). The process control optimization stage focuses on the question, how much binder should be used based on the properties of the base soil material, e.g., how much water should be added to balance the ratio and, if necessary, what is the level of contamination as environmental process parameters.

\subsection{Optimization of binder mixture}
For the two binders, the optimization process is a relatively simple and straightforward process. In this case, all the possible binder combinations can be placed on a straight line that can be described according to equation 1, Figure 1. As long as x or y < 1, the mixture consists of two parts. If x or y=1, the mixture consists of a single binder component. 

\begin{equation}
y = 1 - x
\end{equation}

Possible binder combinations with two different binders are represented by a straight line where the endpoints signify one binder, 1.0 = 100\%, i.e., a one-component mixture. Respectively, adding the \nth{2} binder automatically reduces the amount of the \nth{1} binder, so that together they always constitute a total amount of 100\%, see Figure 1. 
However, for mixtures consisting of three or more binders, the situation becomes more complicated, as the optimal final mixture is hidden among the infinite number of possible mixture combinations. The experimental planning handles this case using simplex method. This method is developed in earlier statistical works (Myers and Montgomery, 1995; Montgomery, 1996) with a general aim to evaluate different types of mixtures with three or more constituent components, which is a very common optimization scenario in real case situations in industry. Figure 2 shows a graph representing all possible mixing ratios for the three components A, B and C, based on the simplex method. The corners represent 100\% of a single binder. The border lines represent mixtures between two different binders, equation (2)-(4). 

\begin{equation}
y = 1 - x	
\end{equation}
\begin{equation}
z = 1 - y
\end{equation}
\begin{equation}
x = 1 - z
\end{equation}
The plane bounded by the edge lines consists of infinitely many combinations of all three binders which can be experimentally mixed and tested. All the possible binder combinations with the three different binders (A, B and C) are represented by a plane bounded by the three straight lines (Figure 2). Here each vertex in the plane represents a single binder, i.e., 1.0 = 100\% of a single component (Myers and Montgomery, 1995).

Depending on the scale and amount of works, chosen for simplex method, both linear effects and interaction effects can be identified. These include, for instance, an interaction effect which occurs when two binder components together produce a significantly higher or lower effect on soil strength than their mutual proportions indicate in advance. The level of details however, places different demands on the total number of tested trials. Therefore, the amount of experimental trails should always be reasonably high. The optimal combination of binders is required for the effective soil stabilisation in geotechnical and economic aspects.

\subsection{Process control optimization}
The process control optimization examines how much binder must be added based on the properties of the base material (e.g., water ratio, organic carbon content or impurity content, the process parameters of soil). In the experimental planning, this optimization was carried out using the factorial experiment. Figure 3 shows a simple 2*2-factor experiment, where two factors are tested at the two levels of various binders (2*2). This method is used to evaluate how two different factors affect the final result. The factor test can be performed with a higher resolution, e.g., by increasing the number of levels to three. In this case, the 2*3-factor trial is then obtained. Similar to the simplex method, higher resolution means that the number of trial experiments increases. The choice of the resolution basically depends on the real project situation (amount, type and quality of soil that should be stabilized) and which interactions are expected to occur in soil treatment (Lindh, 2004).

\subsection{Total binder optimization}
In a total optimization stage, simplex experimental trials from the mixture optimization are combined in a factorial trials with process variables. In this case, the combined experimental set-up then becomes advanced, which is illustrated schematically in Figure 4. The advantage of this approach is that we find an optimum value both in terms of binder combination and binder quantity, based on the important process parameters. Otherwise, if we only investigate the optimum of stabilizing mixture, based on only binder combination, we may miss the unforeseen and unwanted influence from the external process parameters. Of course, the reverse applies if we only optimize binder amount based on the process parameters. This type of the total optimization gives us a robust basis for being able to vary both the binder combinations and the amount of binder in a project, without risking either final quality or project economics. As a result, the combined methodology adopted from Smith (2005) provides both an optimal binder mix and an optimal binder quantity in relation to the process parameters.

%%%%%%%%%%%%%%%%%%%%%%%%%%%%%%%%%%%%%%%%%%
\section{Materials and Methods}

%%%%%%%%%%%%%%%%%%%%%%%%%%%%%%%%%%%%%%%%%%
\section{Results}

%%%%%%%%%%%%%%%%%%%%%%%%%%%%%%%%%%%%%%%%%%
\section{Discussion}

%%%%%%%%%%%%%%%%%%%%%%%%%%%%%%%%%%%%%%%%%%
\section{Conclusions}

%%%%%%%%%%%%%%%%%%%%%%%%%%%%%%%%%%%%%%%%%%
\vspace{6pt} 

%%%%%%%%%%%%%%%%%%%%%%%%%%%%%%%%%%%%%%%%%%
\authorcontributions{For research articles with several authors, a short paragraph specifying their individual contributions must be provided. Conceptualization, methodology, software, validation, visualization, supervision -- Per Lindh; formal analysis, investigation, resources, data curation, project administration, funding acquisition -- Polina Lemenkova; writing---original draft preparation, Polina Lemenkova ; writing---review and editing, Per Lindh and Polina Lemenkova. All authors have read and agreed to the published version of the manuscript.}

\funding{This research was funded by Materials Editorial Office, MDPI, providing 100\% fee waiver for APC for publication of this manuscript.}

\institutionalreview{Not applicable}

\informedconsent{Not applicable}

\dataavailability{Not applicable} 

\acknowledgments{In this section you can acknowledge any support given which is not covered by the author contribution or funding sections. This may include administrative and technical support, or donations in kind (e.g., materials used for experiments).}

\conflictsofinterest{Declare conflicts of interest or state ``The authors declare no conflict of interest.'' Authors must identify and declare any personal circumstances or interest that may be perceived as inappropriately influencing the representation or interpretation of reported research results. Any role of the funders in the design of the study; in the collection, analyses or interpretation of data; in the writing of the manuscript; or in the decision to publish the results must be declared in this section. If there is no role, please state ``The funders had no role in the design of the study; in the collection, analyses, or interpretation of data; in the writing of the manuscript; or in the decision to publish the~results''.} 

\abbreviations{Abbreviations}{
The following abbreviations are used in this manuscript:\\

\noindent 
\begin{tabular}{@{}ll}
MDPI & Multidisciplinary Digital Publishing Institute\\
DOAJ & Directory of open access journals\\
TLA & Three letter acronym\\
LD & Linear dichroism
\end{tabular}
}

%%%%%%%%%%%%%%%%%%%%%%%%%%%%%%%%%%%%%%%%%%
%% Optional
\appendixtitles{no} % Leave argument "no" if all appendix headings stay EMPTY (then no dot is printed after "Appendix A"). If the appendix sections contain a heading then change the argument to "yes".
\appendixstart
\appendix
\section[\appendixname~\thesection]{}
\subsection[\appendixname~\thesubsection]{}
The appendix is an optional section that can contain details and data supplemental to the main text---for example, explanations of experimental details that would disrupt the flow of the main text but nonetheless remain crucial to understanding and reproducing the research shown; figures of replicates for experiments of which representative data are shown in the main text can be added here if brief, or as Supplementary Data. Mathematical proofs of results not central to the paper can be added as an appendix.

\begin{table}[H] 
\caption{This is a table caption.\label{tab5}}
\newcolumntype{C}{>{\centering\arraybackslash}X}
\begin{tabularx}{\textwidth}{CCC}
\toprule
\textbf{Title 1}	& \textbf{Title 2}	& \textbf{Title 3}\\
\midrule
Entry 1		& Data			& Data\\
Entry 2		& Data			& Data\\
\bottomrule
\end{tabularx}
\end{table}

\section[\appendixname~\thesection]{}
All appendix sections must be cited in the main text. In the appendices, Figures, Tables, etc. should be labeled, starting with ``A''---e.g., Figure A1, Figure A2, etc.

%%%%%%%%%%%%%%%%%%%%%%%%%%%%%%%%%%%%%%%%%%
\begin{adjustwidth}{-\extralength}{0cm}
%\printendnotes[custom] % Un-comment to print a list of endnotes

\reftitle{References}

% Please provide either the correct journal abbreviation (e.g. according to the “List of Title Word Abbreviations” http://www.issn.org/services/online-services/access-to-the-ltwa/) or the full name of the journal.
% Citations and References in Supplementary files are permitted provided that they also appear in the reference list here. 

%=====================================
% References, variant A: external bibliography
%=====================================
%\bibliography{your_external_BibTeX_file}

%=====================================
% References, variant B: internal bibliography
%=====================================
\begin{thebibliography}{999}
% Reference 1
\bibitem[Author1(year)]{ref-journal}
Author~1, T. The title of the cited article. {\em Journal Abbreviation} {\bf 2008}, {\em 10}, 142--149.
% Reference 2
\bibitem[Author2(year)]{ref-book1}
Author~2, L. The title of the cited contribution. In {\em The Book Title}; Editor 1, F., Editor 2, A., Eds.; Publishing House: City, Country, 2007; pp. 32--58.
% Reference 3
\bibitem[Author3(year)]{ref-book2}
Author 1, A.; Author 2, B. \textit{Book Title}, 3rd ed.; Publisher: Publisher Location, Country, 2008; pp. 154--196.
% Reference 4
\bibitem[Author4(year)]{ref-unpublish}
Author 1, A.B.; Author 2, C. Title of Unpublished Work. \textit{Abbreviated Journal Name} year, \textit{phrase indicating stage of publication (submitted; accepted; in press)}.
% Reference 5
\bibitem[Author5(year)]{ref-communication}
Author 1, A.B. (University, City, State, Country); Author 2, C. (Institute, City, State, Country). Personal communication, 2012.
% Reference 6
\bibitem[Author6(year)]{ref-proceeding}
Author 1, A.B.; Author 2, C.D.; Author 3, E.F. Title of presentation. In Proceedings of the Name of the Conference, Location of Conference, Country, Date of Conference (Day Month Year); Abstract Number (optional), Pagination (optional).
% Reference 7
\bibitem[Author7(year)]{ref-thesis}
Author 1, A.B. Title of Thesis. Level of Thesis, Degree-Granting University, Location of University, Date of Completion.
% Reference 8
\bibitem[Author8(year)]{ref-url}
Title of Site. Available online: URL (accessed on Day Month Year).
\end{thebibliography}

% If authors have biography, please use the format below
%\section*{Short Biography of Authors}
%\bio
%{\raisebox{-0.35cm}{\includegraphics[width=3.5cm,height=5.3cm,clip,keepaspectratio]{Definitions/author1.pdf}}}
%{\textbf{Firstname Lastname} Biography of first author}
%
%\bio
%{\raisebox{-0.35cm}{\includegraphics[width=3.5cm,height=5.3cm,clip,keepaspectratio]{Definitions/author2.jpg}}}
%{\textbf{Firstname Lastname} Biography of second author}

% For the MDPI journals use author-date citation, please follow the formatting guidelines on http://www.mdpi.com/authors/references
% To cite two works by the same author: \citeauthor{ref-journal-1a} (\citeyear{ref-journal-1a}, \citeyear{ref-journal-1b}). This produces: Whittaker (1967, 1975)
% To cite two works by the same author with specific pages: \citeauthor{ref-journal-3a} (\citeyear{ref-journal-3a}, p. 328; \citeyear{ref-journal-3b}, p.475). This produces: Wong (1999, p. 328; 2000, p. 475)

%%%%%%%%%%%%%%%%%%%%%%%%%%%%%%%%%%%%%%%%%%
%% for journal Sci
%\reviewreports{\\
%Reviewer 1 comments and authors’ response\\
%Reviewer 2 comments and authors’ response\\
%Reviewer 3 comments and authors’ response
%}
%%%%%%%%%%%%%%%%%%%%%%%%%%%%%%%%%%%%%%%%%%
\end{adjustwidth}
\end{document}

